\section{Certified moments from the original arithmetic theta source}
\label{sec:arithmetic-moments}

The moment matrices in the arithmetic kernel formulas can be calculated
directly from the specified theta source.  This section gives an exact
incomplete-gamma series, proves an explicit bound for the entire omitted
double tail, and applies directed interval arithmetic to the retained
terms.  The construction retains the original factor \(g=2\xi\), the
measure \(dx\), the Euler operator \(D=-x\partial_x\), and every
polynomial coefficient.  No information about the location of a zeta
zero enters this calculation.

\subsection{Source polynomials and the exact reflection map}

Keep
\[
 \phi_*(x)=(4\pi^2x^4-6\pi x^2)e^{-\pi x^2},\qquad
 f_j=D^j\Theta\phi_*,\qquad
 \Theta\phi(x)=2\sum_{n\geq1}\phi(nx).
 \tag{AM.1}
\]
Introduce polynomials with their original integer coefficients by
\[
 P_0(z)=4z^2-6z,\qquad
 P_{j+1}(z)=2z\bigl(P_j(z)-P'_j(z)\bigr),\qquad
 Q_j(z)=\sum_{k=0}^j(-1)^k\binom jkP_k(z).
 \tag{AM.2}
\]
Write \(p_{ja}\) and \(q_{ja}\) for the coefficients of \(z^a\).
Both polynomials have degree \(d_j=j+2\), with leading coefficients
\(2^{j+2}\) and \((-1)^j2^{j+2}\), respectively, and both have zero
constant coefficient.  The formula defining \(Q_j\) is an explicit
change of polynomial coordinates; neither family is discarded.

Define the linear involution on the actual Hilbert observation by
\[
 \mathcal T:L^2(\mathbb R_{>0},dx)\longrightarrow
 L^2(\mathbb R_{>0},dx),\qquad
 (\mathcal TF)(x)=x^{-1}F(1/x).
 \tag{AM.3}
\]
It is unitary and self-adjoint, satisfies \(\mathcal T^2=1\), and,
on the strong Schwartz space \(\mathscr B\),
\[
 D\mathcal T=\mathcal T(1-D),\qquad
 \mathcal Tf_0=f_0,\qquad
 \mathcal Tf_j=(1-D)^jf_0.
 \tag{AM.4}
\]
Consequently, the complete source series and its reflected series are
\[
 f_j(x)=2\sum_{n\geq1}P_j(\pi n^2x^2)e^{-\pi n^2x^2},\qquad
 \widetilde f_j(x):=(1-D)^jf_0(x)
   =2\sum_{n\geq1}Q_j(\pi n^2x^2)e^{-\pi n^2x^2},
 \tag{AM.5}
\]
and \(f_j(1/x)=x\widetilde f_j(x)\).

\begin{proof}
For \(z=\pi x^2\), applying \(-x\partial_x\) to
\(P(z)e^{-z}\) gives \(2z(P-P')e^{-z}\).  This proves the
recurrence and (AM.5).  The leading coefficients and zero constant
coefficients follow by induction.  Under the Fourier convention
\(\widehat\phi(\xi)=\int\phi(x)e^{-2\pi ix\xi}\,dx\),
differentiating the Fourier transform of the Gaussian gives
\[
 \widehat{x^2e^{-\pi x^2}}(\xi)
  =\left(\frac1{2\pi}-\xi^2\right)e^{-\pi\xi^2},\qquad
 \widehat{x^4e^{-\pi x^2}}(\xi)
  =\left(\xi^4-\frac{3\xi^2}{\pi}
                    +\frac3{4\pi^2}\right)e^{-\pi\xi^2}.
\]
Multiplication by the exact coefficients \(4\pi^2,-6\pi\)
therefore gives \(\widehat\phi_* =\phi_*\).  Its value at zero
is zero, and its integral is \(\widehat\phi_*(0)=0\).
The theta--Fourier identity already proved for the source \(V\)
then gives \(\mathcal Tf_0=f_0\).

Substitution \(u=1/x\) gives
\(\int|x^{-1}F(1/x)|^2dx=\int|F(u)|^2du\), so \(\mathcal T\)
is unitary.  It is an involution and hence self-adjoint.  Direct
differentiation proves \(D\mathcal T=\mathcal T(1-D)\), including
the constant \(1\) arising from the original factor \(x^{-1}\).
Induction proves the other identities in (AM.4).  Local uniform
convergence of each differentiated Gaussian series follows on
\(x\geq\delta>0\) by a polynomial times Gaussian majorant.  The
series gives decay of all Euler derivatives at infinity, and its
exact reflected expression gives the corresponding decay at zero.
Thus these operations apply to the original functions in
\(\mathscr B\).
\end{proof}

\subsection{An absolutely convergent double series for every entry}

For integers \(i,j\geq0\), retain the raw moment entry
\[
 H_{ij}=\langle D^if_0,D^jf_0\rangle_{L^2(dx)}
  =\int_{\mathbb R}
    \overline{(1/2+it)^i}(1/2+it)^j
       |2\xi(1/2+it)|^2\,\frac{dt}{2\pi}.
 \tag{AM.6}
\]
The equality is the original Mellin Plancherel identity.  In
particular \(H_{ij}\) is real because the source functions are real.
For \(L\geq0\) and \(\lambda>0\), set
\[
 I_L(\lambda)=\int_1^\infty x^{2L}e^{-\lambda x^2}\,dx
   =\frac{\Gamma(L+1/2,\lambda)}{2\lambda^{L+1/2}},
 \qquad
 \Gamma(a,\lambda)=\int_\lambda^\infty u^{a-1}e^{-u}\,du.
 \tag{AM.7}
\]

\begin{theorem}[Exact original arithmetic moment series]
With \(\lambda_{nm}=\pi(n^2+m^2)\), the following double series
converges absolutely, including after the finite coefficient sums
are replaced by their absolute values:
\[
 \boxed{\begin{aligned}
 H_{ij}={}&2\sum_{n,m\geq1}\sum_{a=1}^{d_i}\sum_{b=1}^{d_j}
 (p_{ia}p_{jb}+q_{ia}q_{jb})\pi^{a+b}n^{2a}m^{2b}\\[-2pt]
 &\hspace{24mm}\cdot
 \frac{\Gamma(a+b+1/2,\lambda_{nm})}
      {\lambda_{nm}^{a+b+1/2}}.
 \end{aligned}}
 \tag{AM.8}
\]
Equivalently, it is the Gram integral on the positive half-domain
\[
 H_{ij}=\int_1^\infty
       \bigl(f_i(x)f_j(x)+\widetilde f_i(x)\widetilde f_j(x)\bigr)dx.
 \tag{AM.9}
\]
Thus the precise map to the positive Gram presentation is
\[
 \mathbb C[s]\longrightarrow L^2([1,\infty),dx)^{\oplus2},
 \qquad P\longmapsto\bigl(P(D)f_0,P(1-D)f_0\bigr)\big|_{[1,\infty)}.
 \tag{AM.10}
\]
It preserves the original polynomial inner product (AM.6) and is
injective.  The signs in (AM.8) are retained; the positivity is
the Gram positivity of this specified map.
\end{theorem}

\begin{proof}
Split the integral in (AM.6) at \(x=1\).  On \((0,1)\), set
\(x=1/u\).  The two factors \(u\) supplied by (AM.5) cancel the
Jacobian \(u^{-2}\) exactly.  This gives (AM.9).
For absolute convergence, write
\[
 C_j=\sum_{a=1}^{d_j}|p_{ja}|\pi^a,\qquad
 E_j=\sum_{a=1}^{d_j}|q_{ja}|\pi^a.
 \tag{AM.11}
\]
For \(n,x\geq1\), the polynomial bounds are
\[
 |P_j(\pi n^2x^2)|\leq C_j n^{2d_j}x^{2d_j},\qquad
 |Q_j(\pi n^2x^2)|\leq E_j n^{2d_j}x^{2d_j}.
\]
Since \(n^2+m^2\geq2\) and \(x^2\geq1\),
\[
 (n^2+m^2)x^2\geq\tfrac12(n^2+m^2)+x^2.
\]
The absolute integrand series is therefore bounded by
\[
 4(C_iC_j+E_iE_j)x^{2(d_i+d_j)}e^{-\pi x^2}
 \left(\sum_{n\geq1}n^{2d_i}e^{-\pi n^2/2}\right)
 \left(\sum_{m\geq1}m^{2d_j}e^{-\pi m^2/2}\right),
\]
which is integrable on \([1,\infty)\).  This proves all required
sum--integral interchanges.  Substitution \(u=\lambda x^2\)
gives (AM.7).  The two theta factors contribute \(4\), and
(AM.7) contributes \(1/2\), leaving exactly the factor \(2\)
in (AM.8).

For complex polynomial coefficients, conjugate the coefficients
of the first factor in the Gram integral.  Equation (AM.9) then
proves the inner-product assertion in (AM.10).  If a nonzero
polynomial \(P\) were in its kernel, its first component would
vanish on \([1,\infty)\).  Put \(U(z)=\sum_jP_j(z)c_j\) when
\(P(s)=\sum_jc_js^j\).  The distinct degrees and nonzero leading
coefficients of the \(P_j\) imply \(U\ne0\).  Multiplying the
vanishing theta series by \(e^{\pi x^2}/2\) gives
\[
 U(\pi x^2)=-\sum_{n\geq2}U(\pi n^2x^2)
                         e^{-\pi(n^2-1)x^2}.
\]
The right side decays exponentially times a fixed polynomial as
\(x\to\infty\), by the same Gaussian bound.  A nonzero polynomial
on the left cannot have that decay.  This contradiction proves
injectivity.
\end{proof}

\subsection{A complete finite truncation certificate}

Let \(H_{ij}^{[N]}\) be (AM.8) with both theta indices restricted
to \(1\leq n,m\leq N\), keeping all coefficients \(p_{ia},q_{ia}\).
For the calculation of this entry choose an integer
\(N\geq\max(d_i,d_j)\).  Define, for \(k=i,j\),
\[
 \begin{aligned}
 r_{k,N}&=\left(\frac{N+2}{N+1}\right)^{2d_k}
                 e^{-\pi(2N+3)},&
 T_{k,N}&=\frac{(N+1)^{2d_k}}{1-r_{k,N}},\\
 A_{k,N}&=\sum_{n=1}^N n^{2d_k}e^{-\pi(n^2-1)}
       +T_{k,N}e^{-\pi((N+1)^2-1)}.
 \end{aligned}
 \tag{AM.12}
\]
All these numbers are specified positive real constants.  The
choice of \(N\) ensures \(r_{k,N}<1\), since
\[
 \log r_{k,N}
  <\frac{2d_k}{N+1}-\pi(2N+3)<2-5\pi<0.
\]

\begin{theorem}[Rigorous complete double-tail bound]
For every \(i,j\geq0\) and \(N\geq\max(i,j)+2\),
\[
 \boxed{\begin{aligned}
 |H_{ij}-H_{ij}^{[N]}|\leq\mathcal E_{ij,N}:={}&
 4(C_iC_j+E_iE_j)
 \bigl(T_{i,N}A_{j,N}+A_{i,N}T_{j,N}\bigr)\\[-2pt]
 &\cdot I_{d_i+d_j}\bigl(\pi((N+1)^2+1)\bigr).
 \end{aligned}}
 \tag{AM.13}
\]
In particular the bound tends to zero as \(N\to\infty\), for each
fixed pair \(i,j\), and provides an explicit finite certificate
for every raw arithmetic moment.
\end{theorem}

\begin{proof}
For \(x\geq1\), the ratio of successive terms of
\(n^{2d_k}e^{-\pi n^2x^2}\) is at most
\[
 \left(\frac{n+1}{n}\right)^{2d_k}e^{-\pi(2n+1)}.
\]
The logarithm of this expression has derivative
\(-2d_k/(n(n+1))-2\pi<0\) for positive real \(n\).
For \(n\geq N+1\) every ratio is consequently at most
\(r_{k,N}\).  Summing the geometric majorant yields
\[
 \sum_{n>N}n^{2d_k}e^{-\pi n^2x^2}
  \leq T_{k,N}e^{-\pi(N+1)^2x^2},\qquad
 \sum_{n\geq1}n^{2d_k}e^{-\pi n^2x^2}
  \leq A_{k,N}e^{-\pi x^2}.
 \tag{AM.14}
\]
For the second inequality, divide by \(e^{-\pi x^2}\), bound
each retained factor \(e^{-\pi(n^2-1)x^2}\) at \(x=1\), and
apply the first inequality to the tail.

Write \(F_{k,N}\) for the finite \(P_k\)-series and
\(\widetilde F_{k,N}\) for the finite \(Q_k\)-series.
Equations (AM.11) and (AM.14) give
\[
 \begin{aligned}
 |f_k(x)|, |F_{k,N}(x)|
  &\leq2C_k x^{2d_k}A_{k,N}e^{-\pi x^2},\\
 |f_k(x)-F_{k,N}(x)|
  &\leq2C_kx^{2d_k}T_{k,N}e^{-\pi(N+1)^2x^2}.
 \end{aligned}
\]
The reflected inequalities have exactly \(E_k\) in place of
\(C_k\).  In each product of (AM.9), use the exact decomposition
\(ab-a_Nb_N=(a-a_N)b+a_N(b-b_N)\).  The resulting bound on the
absolute difference of integrands is
\[
 4(C_iC_j+E_iE_j)
 (T_{i,N}A_{j,N}+A_{i,N}T_{j,N})
 x^{2(d_i+d_j)}e^{-\pi((N+1)^2+1)x^2}.
\]
Its integral is exactly the right side of (AM.13).
Finally, the \(A_{k,N}\) remain bounded for fixed \(k\), the
\(T_{k,N}\) grow at most polynomially, and the integral decays
exponentially in \((N+1)^2\).  For completeness, with
\(L=d_i+d_j\) fixed and \(\lambda\geq2L\), the inequality
\(\log x\leq(x^2-1)/2\) for \(x\geq1\) gives
\[
 I_L(\lambda)\leq e^{-\lambda}
   \int_1^\infty e^{-(\lambda-L)(x^2-1)}dx
 \leq \frac{e^{-\lambda}}{2(\lambda-L)}.
\]
Here the last step uses \(x^2-1\geq2(x-1)\).
This proves the asserted limit with an additional elementary
upper bound when desired.
\end{proof}

The finite matrix \((H_{ij}^{[N]})_{0\leq i,j\leq m}\) is itself
positive definite for every \(N\geq1\).  Indeed it is the Gram
matrix of \((F_{j,N},\widetilde F_{j,N})\) on the same half-domain;
the injectivity proof in (AM.10) applies verbatim to the finite
sum.  This fact does not assign an ordering to its difference
from the full matrix.  The relation to the full matrix is the
entrywise enclosure (AM.13), whose signs and error bounds are
explicit.

\subsection{Exact adjoint, parity, and finite-boundary checks}

On \(\mathscr B\) in \(L^2(dx)\), integration by parts gives
\(\langle DF,H\rangle=\langle F,(1-D)H\rangle\).
Therefore the complete moments satisfy
\[
 H_{i+1,j}+H_{i,j+1}=H_{ij},\qquad
 H_{ij}=\sum_{k=0}^i(-1)^k\binom ikH_{0,j+k},\qquad
 H_{01}=\tfrac12H_{00}.
 \tag{AM.15}
\]
All boundary terms vanish because the original functions have
rapid decay at both ends.  More explicitly,
\[
 (Df_i)f_j+f_i(Df_j)-f_if_j=-\partial_x(xf_if_j),
\]
whose integral on \((0,\infty)\) is zero.  Moving each of the
\(i\) copies of \(D\) in the first argument to the second gives
the middle identity; symmetry gives the last one.

The finite sums retain an exact boundary correction at \(x=1\):
\[
 \begin{aligned}
 H_{i+1,j}^{[N]}+H_{i,j+1}^{[N]}-H_{ij}^{[N]}
 ={}&F_{i,N}(1)F_{j,N}(1)\\[-2pt]
 &-\widetilde F_{i,N}(1)\widetilde F_{j,N}(1).
 \end{aligned}
 \tag{AM.16}
\]
To prove this, use \(DF_{j,N}=F_{j+1,N}\) for the first
component and \(D\widetilde F_{j,N}=\widetilde F_{j,N}
-\widetilde F_{j+1,N}\) for the second.  Applying the preceding
integration-by-parts calculation on \([1,\infty)\) gives the
first displayed endpoint product for the first component and
its negative for the second.  The endpoints at infinity vanish
term by term.  In particular the exact finite identity
\(H_{01}^{[N]}=H_{00}^{[N]}/2\) follows even before the omitted
tail is enclosed, since \(F_{0,N}=\widetilde F_{0,N}\).

One may also retain the exact triangular coordinate map
\[
 B_{ja}=\begin{cases}\binom ja(-1/2)^{j-a},&a\leq j,\\0,&a>j,
             \end{cases}
 \qquad C=BHB^T.
 \tag{AM.17}
\]
This is the Gram matrix of \(u_j=(D-1/2)^jf_0\), with inverse
coordinate matrix \((B^{-1})_{ja}=\binom ja(1/2)^{j-a}\).
No original moment is lost.  Since
\(\mathcal Tu_j=(-1)^ju_j\), unitarity proves \(C_{ij}=0\)
for odd \(i+j\).  The operator \(D-1/2\) is skew-symmetric on
this displayed strong Schwartz domain, so for \(i+j=2q\),
\[
 C_{ij}=(-1)^{i+q}\|u_q\|^2,\qquad
 \|u_q\|^2=\int_{\mathbb R}t^{2q}|2\xi(1/2+it)|^2
                                      \frac{dt}{2\pi}>0.
 \tag{AM.18}
\]
For the first equality move all \(i\) powers to the second
argument, and compare with moving the \(q\) powers in
\(\langle u_q,u_q\rangle\).  Mellin Plancherel gives the
second equality.  Strict positivity follows from the nonzero
original function and the polynomial injectivity already proved.

\subsection{Reproducible certified finite arithmetic data}

The companion program \texttt{scripts/check\_arithmetic\_moments.py}
implements the exact integer recurrence (AM.2), evaluates every
retained term of (AM.8) by Arb directed ball arithmetic, and
enlarges its radius by an upward-rounded value of (AM.13).
Thus the final balls enclose the infinite-source entries, including
rounding error and both omitted theta tails.  The default run
uses \(0\leq i,j\leq5\), \(N=7\), and \(110\) decimal working
digits; its machine-readable output is
\texttt{checks/arithmetic\_moments.json}.

Let \(\Delta_k=\det(H_{ij})_{0\leq i,j<k}\) and
\(\Delta_0=1\).  The squared norm \(h_k\) of the monic
degree-\(k\) polynomial for the original measure \(\mu\) is
\[
 h_k=\frac{\Delta_{k+1}}{\Delta_k}.
 \tag{AM.19}
\]
Indeed the coefficients below the leading term solve the Gram
system making that monic polynomial orthogonal to degrees below
\(k\).  Substitution of that unique solution into its squared
norm gives the Schur complement of the leading \(k\)-by-\(k\)
Gram block.  The block determinant identity gives (AM.19).
The interval calculation certifies strict positivity of all six
leading principal determinants recorded in the companion certificate
and all six corresponding
\(h_k\); those are the actual arithmetic norms used by the
orthogonal recurrence, through the finite degrees stated.

Here are explicit certified intervals from that run.  Each row
means the entire interval \([c-r,c+r]\), with the printed decimal
centre \(c\) and radius \(r\) interpreted exactly.  Rounding the
stored Arb balls to the displayed decimals was performed outwards.
\[
\begin{array}{c|r|l}
 \text{raw moment}&c&r\\\hline
 H_{00}&1.279007247846485140480&4.67\cdot10^{-22}\\
 H_{01}&0.6395036239232425702398&3.33\cdot10^{-23}\\
 H_{02}&-12.73579749060893715027&2.81\cdot10^{-21}\\
 H_{11}&13.37530111453217972051&2.57\cdot10^{-21}\\
 H_{12}&6.687650557266089860256&2.85\cdot10^{-22}\\
 H_{22}&389.8820543216399270934&1.65\cdot10^{-20}
\end{array}
\]
Symmetry supplies the other entries of this original three-by-three
moment block.  The six original monic squared norms are enclosed by
\[
\begin{array}{c|r|l}
 k&c\text{ for }h_k&r\\\hline
 0&1.279007247846485140480&4.67\cdot10^{-22}\\
 1&13.05554930257055843539&2.69\cdot10^{-21}\\
 2&250.0089768190499145175&4.88\cdot10^{-20}\\
 3&6694.666685525627506773&4.76\cdot10^{-19}\\
 4&224298.1953314520556463&9.95\cdot10^{-18}\\
 5&9383436.471794727053469&2.54\cdot10^{-16}
\end{array}
\]
The analytic double-tail bound is less than
\(2.961369\cdot10^{-63}\) for every entry of the six-by-six
matrix, and is less than \(1.150030\cdot10^{-83}\) for \(H_{00}\).
These bounds account for omitted source terms; the displayed
intervals additionally include finite-term arithmetic and decimal
conversion errors.

The checker also evaluates (AM.15) and the independent exact
finite-boundary identity (AM.16), and detects removal of either
the reflected half-domain, the theta factor \(2\), or the
adjoint constant \(1\).  Its checks execute identically under
ordinary Python and Python's optimization flag.  Finite interval
certificates do not assert an infinite-degree weight estimate;
(AM.8)--(AM.14) give the exact arithmetic input and its complete
error control for calculating further degrees.
